\title{Coordinated learning}
\author{Luciano Venturim}
% \author{Luciano Venturim\footnote{Northwestern University Kellogg School of Management, Global Hub, 2211 Campus Drive, Evanston IL 60208; email: lucianofabio.venturim@kellogg.northwestern.edu.}}
\date{\today}
\documentclass[12pt,letterpaper]{article}

% =========================
% Encoding & fonts
% =========================
\usepackage[T1]{fontenc}
\usepackage[utf8]{inputenc}  % UTF-8 source files (Overleaf default)
% \usepackage{newtxtext}       % Times-like text
% \usepackage{newtxmath}       % Times-like math
\usepackage{microtype}       % better kerning & protrusion

% =========================
% Math & theorem tools
% =========================
\usepackage{amsmath,amssymb,amsthm}
\usepackage{bm}              % bold math
\usepackage{mathrsfs}        % \mathscr
\usepackage{bbm}             % \mathbbm{1}
\usepackage{dsfont}          % \mathds

% =========================
% Layout & spacing
% =========================
\usepackage[top=1in, bottom=1in, left=1in, right=1in]{geometry}
\usepackage{setspace}        % \doublespacing, \onehalfspacing, etc.
\usepackage[shortlabels]{enumitem}        % compact lists, label control
\usepackage{float}           % [H] placement
\sloppy                      % relax line-breaking to avoid overfull lines

% =========================
% Graphics, tables, colors
% =========================
\usepackage{graphicx}
\usepackage{tikz}
\usetikzlibrary{patterns}
\usepackage{rotating}
\usepackage{xcolor}          % (supersedes 'color')
\usepackage{booktabs}        % \toprule, \midrule, \bottomrule
\usepackage{array}
\usepackage{multirow}

% =========================
% Text utilities & domain packages
% =========================
\usepackage{csquotes}
\usepackage{comment}
\usepackage{sgame}           % strategic-form games (if used)
\usepackage{econometrics}    % keep only if your project includes this .sty

% =========================
% Bibliography & hyperlinks
% =========================
\usepackage[round]{natbib}
\bibliographystyle{abbrvnat}

\usepackage[pagebackref]{hyperref}  % load hyperref close to last
\hypersetup{
  colorlinks=true,
  linkcolor=black,
  citecolor=black,
  urlcolor=black
}

% =========================
% New commands
% =========================
\newcommand{\argmax}[1]{\underset{#1}{\text{argmax }}}
\newcommand{\argmin}[1]{\underset{#1}{\text{argmin }}}
\newcommand{\R}{\mathbb{R}}
\newcommand{\N}{\mathbb{N}}
\newcommand{\Z}{\mathbb{Z}}
\newcommand{\Q}{\mathbb{Q}}
\newcommand{\Ex}{\mathbb{E}}

\newtheorem{proposition}{\hskip\parindent\bf{Proposition}}
\newtheorem{theorem}{\hskip\parindent\bf{Theorem}}
\newtheorem{lemma}{\hskip\parindent\bf{Lemma}}
\newtheorem{corollary}{\hskip\parindent\bf{Corollary}}
\newtheorem{condition}{\hskip\parindent\bf{Condition}}
\newtheorem{claim}{\hskip\parindent\bf{Claim}}
\newtheorem{definition}{\hskip\parindent\bf{Definition}}
\newtheorem{assumption}{\hskip\parindent\bf{Assumption}}
\newtheorem{example}{\hskip\parindent\bf{Example}}

\begin{document}

\doublespacing

\maketitle

% \begin{abstract}
% ABSTRACT. \\
% JEL Codes:. Keywords.
% \end{abstract}

% \newpage
\section{Model}

Consider the following model. A continuum of unit mass of individuals $i\in[0,1]$ play a game repeatedly. At every period, two individuals are matched, with matches being formed at random. Within that period, the pair of individuals $i,j$ play a simultaneous stage game $G(\theta)$. The payoffs obtained in a stage game vary with the unknown state of the world $\theta\in\Theta$. After period $t$'s actions, each individual survives with probability $\delta\in(0,1)$. The individuals who died are replaced by new individuals with a clean history. Individuals' payoffs are normalized by $(1-\delta)$. We also allow the individuals to report their history of play and payoffs in a credible and \textbf{costless} manner. The communication happens after the match is realized, but before the actions are taken.

\paragraph{Example.} For now, let's consider the case in which $\theta$ is persistent --- it is drawn at $t=0$ according to prior distribution $p\in \Delta(\Theta)$ and it stays the same from then on ---, and there are no replacements.\footnote{In this setting, we can either interpret $\delta$ as a discount factor, or continue assuming that agents are replaced but that new individuals inherit the histories.} Consider the following coordination example. There are two states indexed by $r\in\{r_1,r_2\}$. Players can play $a_i\in\{L,R\}$. The payoffs are given by 

\begin{table}[H]
    \centering
    \begin{tabular}{|c|c|c|}
        \hline
        $i\,\backslash \, j$  & L & R\\
        \hline
        L & $2,2$ & $3,-1$\\
        \hline
        R & $-1,3$ & $r,r$\\
        \hline
    \end{tabular}
\end{table}

Note that only action profile $a=(R,R)$ differentiates the state. Suppose $r_1$, $r_2$, and $p_1=P(r=r_1)$ are such that $$r_2>3>2>r_1$$ and $$3>p_1r_1+(1-p_1)r_2>2.$$ Under this assumption, the following are true:

\begin{enumerate}[i.]
    \item When $r=r_1$, $a_i=L$ is a strictly dominating action. Hence, $(L,L)$ is the unique strict, Pareto Efficient Nash Equilibrium (NE) of the stage game.

    \item When $r=r_2$, both $(L,L)$ and $(R,R)$ are strict NE, with $(R,R)$ Pareto dominating $(L,L)$.

    \item For all $\delta\in(0,1)$, the third condition implies that it is efficient for everyone to play $(R,R)$ at $t=0$ and coordinate in the efficient equilibrium thereafter. This can be clearly seen by noting that \begin{equation*}
        p_1[(1-\delta)r_1+2\delta]+(1-p_1)r_2>p_1r_1+(1-p_1)r_2>2.
    \end{equation*}

    \item When $\delta$ is low, deviating at $t=0$ is profitable even if leads to the deviating player not learning the state.
\end{enumerate}

In this efficient outcome, everyone tries the ``risky" action at $t=0$, learn the state and play the efficient action from them one. Can it be an equilibrium? It clearly depends on what is observable when a new match is formed. If no record can be tracked, one can show that this outcome can not be an equilibrium for some parameters for all discount factors.

First, note that a deviation from the strategy at any $t\geq 1$ is clearly not profitable. Consider then a deviation from $a_{i,0}=R$ to $\hat{a}_{i,0}=L$. Since $u_i(\hat{a}_{i,0},R,\theta_1)=u_i(\hat{a}_{i,0},R,\theta_2)$, player $i$ does not learn the state and hence his posterior equals his prior $p$. At $t=1$, in his new match, the other player will play $L$ with probability $p_1$ and $R$ with probability $1-p_1$. In any case, player $i$ can learn the state at $t=1$ from the action of his match.

Now, consider the case $2p_1+3(1-p_1)>-p_1+r_2(1-p_1)$, so that player $i$'s best response at $t=1$ is $L$. The proposed strategy is an equilibrium if and only if \begin{equation*}
    p_1[(1-\delta)r_1+2\delta]+(1-p_1)r_2 \geq 3(1-\delta)+\delta[2p_1+(1-p_1)(3(1-\delta)+\delta r_2)].
\end{equation*}

Note that both the LHS and RHS converge to $2p_1+(1-p_1)r_2$ as $\delta$ converges to $1$. On the other hand, the LHS converges to $p_1r_1+(1-p_1)r_2$ when $\delta \to 0$, while the RHS converges to $3$. Finally, since the LHS is linear in $\delta$ and the RHS is quadratic, there exists at most a discount factor between $0$ and $1$ for which this is an equilibrium. In fact, one can check that if $r_1=1$, $r_2=4$, and $p_1=1/2$, this inequality is only satisfied at $\delta=1$. This is also true for $r_2\in[3,4]$. In those cases, the proposed strategy is not an equilibrium for any $\delta$ below unity.

\begin{lemma}
    In the game above, when $r_1=1$, $r_2\in[3,4]$, and $p_1=1/2$, full-learning and coordination is not an equilibrium for any $\delta<1$.
\end{lemma}

Now, consider the case in which $r_1=1$, $r_2=4$, and $p_1=1/2$, but suppose the players have access to a public randomization device at $t=0$, which allows coordination at different strategies for different matches. Consider a strategy profile in which $\varepsilon>0$ of the players at $t=0$ (or $\frac{\varepsilon}{2}$ of the matches) are assigned to play $(L,L)$ at $t=0$, while the others are assigned to play $(R,R)$. At $t=1$ onward, players who have played $R$ at $t=0$ coordinate on the efficient equilibrium, while players who have played $L$ play $L$ at all histories in which non of their past partners have played $R$, and $R$ at all histories in which at least one of their partners have played $R$. Note that on the equilibrium path, $a_{-i}=R$ can only be observed if $r=r_2$. Hence, observing $a_{-i}=R$ signals that your partner learned the state at a prior date, and therefore player $i$ updates his posterior to attach probability $1$ to $r=r_2$. On the other hand, player $i$ can observe $a_{-i}=L$ by either matching with an uninformed player or by matching with an informed player who learned that the state is $r=r_1$. After observing $a_{-i,1}=L$, the player attaches probability close to $1$ to the second case when $\varepsilon$ is sufficiently small, in which case playing $L$ is a best response. The posterior probability that player $i$ attaches to state $r=r_1$ increases over time with more observations of $a_{-i}=L$, and changes to $0$ after any observation of $a_{-i}=R$.

% \newpage
% %\nocite{*}
% \bibliography{references} % Entries are in the refs.bib file

% \newpage
% \appendix

% \section{Appendix}

% %Restating Propositions in the text with the same number
% \begingroup
% \renewcommand{\theproposition}{\ref{REFERENCE}}
% \begin{proposition}[Restated]
%     PROPOSITION

%     \begin{proof}
%         PROOF.
%     \end{proof}
% \end{proposition}
% \endgroup

\end{document}
